% GENERATED FILE. Edit canonical lessons, then run npm run generate:books.
% canonical-source-hash: fnv1a64:0faf3c53a88f95a6
% canonical-lessons: JA-C265-kemuri, JA-C265-hirogaru, JA-C265-shoubousha, JA-C265-nigemichi, JA-C265-kaketsukeru

\chapter{Smoke on the Stairs}
\label{ch:ja-smoke-on-the-stairs}
\hlchaptermodality{265}

\begin{chapteropening}
\textbf{By the end of this chapter:} \emph{I can say what happened in a fire and who came.}

Complete the last lesson of chapter 265: I can say what happened in a fire and who came.
\end{chapteropening}

\section[\texorpdfstring{\ja{けむり} (kemuri)}{けむり (kemuri)}]{\ja{けむり} (kemuri) — smoke}

\label{lesson:JA-C265-kemuri}



\begin{quote}
Before the new one: say the Japanese for very, then the Japanese for for a while.
\end{quote}

\subsection*{You'll want to know: \ja{けむり}}
\textbf{\ja{けむり}} — \emph{kemuri} — \textquotedblleft{}smoke\textquotedblright{}.

\begin{quote}
\textbf{\ja{けむりで} \ja{そらが} \ja{くらく} \ja{なりました。}} — \emph{kemuri de sora ga kuraku narimashita} — \textquotedblleft{}The sky went dark with smoke.\textquotedblright{}
\end{quote}

\begin{cousinweb}
\textbf{\ja{けむり}} (\emph{kemuri}) is a native Japanese word (wago), inherited rather than borrowed from Chinese; its kanji was attached later as a kun'yomi (JMdict).
\end{cousinweb}

\begin{grammarlens}[title={What you've built}]
One more word for telling what happened.
\end{grammarlens}

\subsection*{Your turn}
\begin{itemize}
\raggedright
  \item \emph{Say it:} \emph{kemuri}
  \item \emph{Say it:} \emph{kemuri}, once more
  \item \emph{Say it:} \emph{kemuri de sora ga kuraku narimashita} — The sky went dark with smoke.
  \item \emph{Recall:} say the Japanese for very, then the Japanese for for a while, then say \emph{kemuri} again
\end{itemize}

\begin{tcolorbox}[breakable,colback=teal!4,colframe=teal!35!black,title={Before you move on}]
What does \ja{けむり} mean? (\textquotedblleft{}Smoke\textquotedblright{}.) Say it once more.
\end{tcolorbox}
\section[\texorpdfstring{\ja{ひろがる} (hirogaru)}{ひろがる (hirogaru)}]{\ja{ひろがる} (hirogaru) — to spread; to widen}

\label{lesson:JA-C265-hirogaru}



\begin{quote}
Before the new one: say the Japanese for for a while, then the Japanese for smoke.
\end{quote}

\subsection*{You'll want to know: \ja{ひろがる}}
\textbf{\ja{ひろがる}} — \emph{hirogaru} — \textquotedblleft{}to spread; to widen\textquotedblright{}.

\begin{quote}
\textbf{\ja{かぜが} \ja{あったら、ひが} \ja{ひろがるでしょう。}} — \emph{kaze ga attara, hi ga hirogaru deshō} — \textquotedblleft{}If there were a wind, the fire would spread.\textquotedblright{}
\end{quote}

Polite \textbf{\ja{ひろがります}}, \emph{hirogarimasu}; a godan verb. Past (polite): \textbf{\ja{ひが} \ja{ひろがりました。}}, \emph{hi ga hirogarimashita}.

\begin{cousinweb}
\textbf{\ja{ひろがる}} (\emph{hirogaru}) is a native Japanese word (wago), inherited rather than borrowed from Chinese; its kanji was attached later as a kun'yomi (JMdict). It is the intransitive partner of \emph{hirogeru}, \textquotedblleft{}to spread out\textquotedblright{}.
\end{cousinweb}

\begin{grammarlens}[title={What you've built}]
One more word for telling what happened.
\end{grammarlens}

\subsection*{Your turn}
\begin{itemize}
\raggedright
  \item \emph{Say it:} \emph{hirogaru}
  \item \emph{Say it:} \emph{hirogaru}, once more
  \item \emph{Say it:} \emph{kaze ga attara, hi ga hirogaru deshō} — If there were a wind, the fire would spread.
  \item \emph{Recall:} say the Japanese for for a while, then the Japanese for smoke, then say \emph{hirogaru} again
\end{itemize}

\begin{tcolorbox}[breakable,colback=teal!4,colframe=teal!35!black,title={Before you move on}]
What does \ja{ひろがる} mean? (\textquotedblleft{}To spread; to widen\textquotedblright{}.) Say it once more.
\end{tcolorbox}
\section[\texorpdfstring{\ja{しょうぼうしゃ} (shōbōsha)}{しょうぼうしゃ (shōbōsha)}]{\ja{しょうぼうしゃ} (shōbōsha) — a fire engine}

\label{lesson:JA-C265-shoubousha}



\begin{quote}
Before the new one: say the Japanese for smoke, then the Japanese for to spread.
\end{quote}

\subsection*{You'll want to know: \ja{しょうぼうしゃ}}
\textbf{\ja{しょうぼうしゃ}} — \emph{shōbōsha} — \textquotedblleft{}a fire engine\textquotedblright{}.

\begin{quote}
\textbf{\ja{しょうぼうしゃが} \ja{すぐ} \ja{きました。}} — \emph{shōbōsha ga sugu kimashita} — \textquotedblleft{}A fire engine arrived at once.\textquotedblright{}
\end{quote}

\begin{cousinweb}
\textbf{\ja{しょうぼうしゃ}} (\emph{shōbōsha}) is Sino-Japanese (kango): three Chinese characters read with their on'yomi, readings borrowed from Chinese. They are \emph{shō}, \textquotedblleft{}extinguish\textquotedblright{} (Middle Chinese \emph{sjew}), \emph{bō}, \textquotedblleft{}ward off\textquotedblright{} (\emph{bjang}) and \emph{sha}, \textquotedblleft{}car\textquotedblright{} (\emph{tsyhae}) (KANJIDIC2; Middle Chinese in Baxter's transcription).
\end{cousinweb}

\begin{grammarlens}[title={What you've built}]
One more word for telling what happened.
\end{grammarlens}

\subsection*{Your turn}
\begin{itemize}
\raggedright
  \item \emph{Say it:} \emph{shōbōsha}
  \item \emph{Say it:} \emph{shōbōsha}, once more
  \item \emph{Say it:} \emph{shōbōsha ga sugu kimashita} — A fire engine arrived at once.
  \item \emph{Recall:} say the Japanese for smoke, then the Japanese for to spread, then say \emph{shōbōsha} again
\end{itemize}

\begin{tcolorbox}[breakable,colback=teal!4,colframe=teal!35!black,title={Before you move on}]
What does \ja{しょうぼうしゃ} mean? (\textquotedblleft{}A fire engine\textquotedblright{}.) Say it once more.
\end{tcolorbox}
\section[\texorpdfstring{\ja{にげみち} (nigemichi)}{にげみち (nigemichi)}]{\ja{にげみち} (nigemichi) — an escape route; a way out}

\label{lesson:JA-C265-nigemichi}



\begin{quote}
Before the new one: say the Japanese for to spread, then the Japanese for a fire engine.
\end{quote}

\subsection*{You'll want to know: \ja{にげみち}}
\textbf{\ja{にげみち}} — \emph{nigemichi} — \textquotedblleft{}an escape route; a way out\textquotedblright{}.

\begin{quote}
\textbf{\ja{にげみちは} \ja{かいだんです。}} — \emph{nigemichi wa kaidan desu} — \textquotedblleft{}The way out is by the stairs.\textquotedblright{}
\end{quote}

\begin{cousinweb}
\textbf{\ja{にげみち}} (\emph{nigemichi}) is a native Japanese word (wago), inherited rather than borrowed from Chinese, built from \emph{nige}, \textquotedblleft{}fleeing\textquotedblright{} (from \emph{nigeru}), and \emph{michi}, \textquotedblleft{}road\textquotedblright{} (JMdict, KANJIDIC2).
\end{cousinweb}

\begin{grammarlens}[title={What you've built}]
One more word for telling what happened.
\end{grammarlens}

\subsection*{Your turn}
\begin{itemize}
\raggedright
  \item \emph{Say it:} \emph{nigemichi}
  \item \emph{Say it:} \emph{nigemichi}, once more
  \item \emph{Say it:} \emph{nigemichi wa kaidan desu} — The way out is by the stairs.
  \item \emph{Recall:} say the Japanese for to spread, then the Japanese for a fire engine, then say \emph{nigemichi} again
\end{itemize}

\begin{tcolorbox}[breakable,colback=teal!4,colframe=teal!35!black,title={Before you move on}]
What does \ja{にげみち} mean? (\textquotedblleft{}An escape route; a way out\textquotedblright{}.) Say it once more.
\end{tcolorbox}
\section[\texorpdfstring{\ja{かけつける} (kaketsukeru)}{かけつける (kaketsukeru)}]{\ja{かけつける} (kaketsukeru) — to rush to the scene}

\label{lesson:JA-C265-kaketsukeru}



\begin{quote}
Before the new one: say the Japanese for a fire engine, then the Japanese for an escape route.
\end{quote}

\subsection*{You'll want to know: \ja{かけつける}}
\textbf{\ja{かけつける}} — \emph{kaketsukeru} — \textquotedblleft{}to rush to the scene\textquotedblright{}.

\begin{quote}
\textbf{\ja{おまわりさんが} \ja{すぐ} \ja{かけつけました。}} — \emph{omawarisan ga sugu kaketsukemashita} — \textquotedblleft{}A police officer was there in no time.\textquotedblright{}
\end{quote}

Polite \textbf{\ja{かけつけます}}, \emph{kaketsukemasu}; an ichidan verb. Past (polite): \textbf{\ja{しょうぼうしゃが} \ja{かけつけました。}}, \emph{shōbōsha ga kaketsukemashita}.

\begin{cousinweb}
\textbf{\ja{かけつける}} (\emph{kaketsukeru}) is a native Japanese word (wago), inherited rather than borrowed from Chinese; its kanji was attached later as a kun'yomi (JMdict).
\end{cousinweb}

\begin{grammarlens}[title={What you've built}]
One more word for telling what happened.
\end{grammarlens}

\subsection*{Your turn}
\begin{itemize}
\raggedright
  \item \emph{Say it:} \emph{kaketsukeru}
  \item \emph{Say it:} \emph{kaketsukeru}, once more
  \item \emph{Say it:} \emph{omawarisan ga sugu kaketsukemashita} — A police officer was there in no time.
  \item \emph{Recall:} say the Japanese for a fire engine, then the Japanese for an escape route, then say \emph{kaketsukeru} again
\end{itemize}

\begin{tcolorbox}[breakable,colback=teal!4,colframe=teal!35!black,title={Before you move on}]
What does \ja{かけつける} mean? (\textquotedblleft{}To rush to the scene\textquotedblright{}.) Say it once more.
\end{tcolorbox}
